\documentclass[lang=cn]{ztex}
\usepackage{zhlipsum}


\setcounter{tocdepth}{6}
\setcounter{secnumdepth}{6}
\ztocenable
%% NOTE: mindmap.lib.tex must load after command '\ztocenable' !!!
\usepackage{tikz}
\usepackage{tikz-qtree}
\usepackage{tikz-qtree-patch}

\definecolor{Red}{HTML}{e32b2d}
\definecolor{Green}{HTML}{15821b}
\definecolor{Blue}{HTML}{0984e3}
\tikzset{
  section/.style={draw, rectangle, rounded corners, fill=Red},
  subsection/.style={draw, rectangle, rounded corners, fill=Blue, text=white},
  subsubsection/.style={draw, rectangle, rounded corners, fill=Green, text=white},
}

\ExplSyntaxOn
\def\toclinecmd#1#2#3#4{{#3}~}
\ztoc_group_parser:NNNnnnnn \g_ztoc_keyvaltoc_seq
  \toclinecmd \g_tmpa_tl {[.}{}{}{~]}{}
\def\TocTree#1
  {
    \exp_args:Nno \__toc_tree:nn
      {#1}{\g_tmpa_tl}
  }
\cs_new:Npn \__toc_tree:nn #1#2
  {
    \quark_if_no_value:nF {#1}
      {
        \Tree[.\node[section]{#1};~#2 ]
      }
  }
\ExplSyntaxOff

\NewDocumentEnvironment{mindmap}{}{%
  \begin{tikzpicture}[
    grow'=right, level distance=1cm,
    sibling distance=10pt,
    every tree node/.style = {
      anchor=west, outer sep=0pt, draw, rectangle, 
      rounded corners, fill=gray!25
    },
    every level 1 node/.style = {
      subsection
    },
    every level 2 node/.style = {
      subsubsection
    },
    edge from parent/.style={
      draw, thick, rounded corners,
      edge from parent path={
        (\tikzparentnode.east) -- +(.5cm, 0pt)
        |- (\tikzchildnode.west)
      },
    },
  ]
}{\end{tikzpicture}}


%% rebustfy some cmds:
\ExplSyntaxOn
\clist_map_function:nN 
  { \vec, \bigcup, \mathbb, \mathrm }
  \zcmd_robustify:N


% %% re-define sectioning commands format
\zsecformat{\section, \subsection, \subsubsection, \paragraph}
  {
    explicit=true, code={
      \exp_args:Nooo \__zsect_title_toc_add:nnn { \zsecclass }{ \zsectocname }{ \zsecname }
    }
  }
\ExplSyntaxOff


\begin{document}
\section{点}
\subsection{内点}
\subsubsection{存在 $\vec{x}$ 的一个 $\delta$ 邻域 $O(\vec{x}, \delta)$ 完全落在 $S$ 中}
\subsubsection{全体的点记作 $\mathrm{int}\, S$ 或 $S^\circ$}


\subsection{界点}
\subsection{外点}
\subsection{聚点}
    \subsubsection{定义}
        \paragraph{定义1： $\vec{x}$ 的任一邻域都包含 $S$ 中无数个点.}
        \paragraph{定义2： $\vec{x}$ 的任一邻域都包含 $S$ 中的至少一个点.}
    \subsubsection{全体的点记作: $S^d$ 或 $S'$}
    \subsubsection{直观理解: 也就是点 $\vec{x}$ 附近是不是聚集了很多 $S$ 中的点.}
    \subsubsection{内点一定是聚点，界点一定是聚点, 有限集没有聚点.}
\subsection{孤立点}




\section{点集}
\subsection{开集}
\subsection{相对开集}
\subsection{闭集}
\subsection{相对闭集}
\subsection{紧集}
    \subsubsection{开覆盖}
        \paragraph{如果 $\mathbb{R}^n$ 的一组开覆盖 $\{U_\alpha\}$ 满足 $\displaystyle\bigcup_{\alpha}U_\alpha \supset S$, 那么称 $\{U_\alpha \}$ 是 $S$ 的一个开覆盖.}
        \paragraph{设 $E$ 是任何一个度量空间的子集, 为每个 $p\in E$ 选择一个半径 $r_p$（半径可以是不同的）,那么集族 $N_{r_p}(P)|p\in E$ 是 $E$ 的一个开覆盖}
    \subsubsection{紧集}
    \subsubsection{Heine-Borel 定理}
        \paragraph{$\mathbb{R}^n$ 上的点集 $S$ 是紧集的充分必要条件是它为有界闭集}
        \paragraph{一个测试插入图片: \lower3em\hbox{\includegraphics[width=10em]{example-image-a}}}


\section{区域}
\subsection{开域}
\subsection{闭域}


\section{基本定理}
    \subsection{闭矩形套定理}
    \subsection{Cantor 闭区间套定理}
    \subsection{凝聚定理(Bolzano-Weierstrass 定理): $\mathbb{R}^n$ 中的有界点列一定有收敛子列(有界无限点集一定有聚点).}
    \subsection{Cauchy 定理:收敛点列 $\leftrightarrow$ 基本点列}
    \subsection{紧性定理(Heine-Borel 定理): 一个基本测试}


% \newpage
% Insert a new mindmap here, which relies on file: \texttt{mindmap.lib.tex}.

% \tableofcontents

% \newpage
\pagestyle{empty}
% \pdfpagewidth 22.5in\relax
% \pdfpageheight 16.75in\relax

\zhlipsum[1]


\begin{figure}[!htbp]
  \centering
  \wscale{.75\linewidth}{%
    \begin{mindmap}
      \TocTree{\parbox{12em}{\centering\textbf{平面点集}\\($S$ 为 $\mathbb{R}^n$ 上的点集, 其补集 $\mathbb{R}^n\backslash S$ 记作 $S^c$)}}
    \end{mindmap}}
  
  \caption{Mindmap Example}
  \label{fig:mind-map}
\end{figure}
\zhlipsum[2]
\end{document}
